% Prova 01 (Macroeconomics I, MPE Insper, 2026/T3) --- detailed solution.
% Diagnosis of the graded answers: prova-01-error-analysis.md (same folder).
% Format (rules/00_estilo_avaliacao.md §4): every equation used later is numbered and cited at each
% step; each question opens with a Notation list; one operation per displayed line.
% Numbering: Q1 uses (G1)-(G4); Q2 keeps the statement's (1)-(4) and continues from (5); Q3 keeps the
% statement's (AD), (AS), (YN), (LF), (YE) and continues the numeric sequence.
% Build from Prova/: pdflatex prova-01-solution.tex (twice); pdffonts prova-01-solution.pdf (all Type 1).
% Every closed form: python prova-01-solution-check.py
\documentclass[11pt,a4paper]{article}
\usepackage[utf8]{inputenc}
\usepackage[T1]{fontenc}
\usepackage{lmodern}   % Type 1 vector fonts (replaces PK bitmaps)
\usepackage{cmap}      % ToUnicode CMap -> copyable text and accents
\usepackage[portuguese,english]{babel}
\usepackage{amsmath,amssymb}
\usepackage[a4paper,margin=2.2cm]{geometry}
\usepackage{enumitem}
\usepackage{fancyhdr}
\usepackage{xcolor}
\usepackage{microtype}
\usepackage{booktabs}
\usepackage{needspace}
\usepackage{tikz}
\usepackage{pgfplots}
\pgfplotsset{compat=1.16}

\pagestyle{fancy}\fancyhf{}
\lhead{\footnotesize Macroeconomics I --- MPE Insper}
\rhead{\footnotesize Prova 01 --- detailed solution}
\cfoot{\footnotesize\thepage}
\renewcommand{\headrulewidth}{0.3pt}
\setlength{\headheight}{14pt}
\setlist{topsep=2pt,itemsep=1pt}
\allowdisplaybreaks

\newcommand{\questao}[2]{\par\bigskip\needspace{6\baselineskip}\noindent
  {\large\bfseries\color{blue!55!black}Question #1 --- #2}\par\noindent\rule{\linewidth}{0.3pt}\par}
% The exam's statement, verbatim and in the exam's language (Portuguese), as printed.
\newcommand{\enunciado}[1]{\par\smallskip\noindent{\footnotesize\textit{Statement, as printed in the exam:}}%
  \begin{quote}\small\begin{otherlanguage}{portuguese}#1\end{otherlanguage}\end{quote}}
\newcommand{\itemq}[1]{\par\medskip\needspace{6\baselineskip}\enunciado{#1}}
\newcommand{\passo}[1]{\par\smallskip\noindent\textbf{Step --- #1.}~}
\newcommand{\intuicao}[1]{\par\smallskip\noindent\textit{Economic intuition.}~#1}
\newcommand{\erro}[1]{\par\smallskip\noindent{\color{red!60!black}\textbf{Where the marks went.}~#1}}
\newcommand{\refbib}[1]{\par\smallskip\noindent{\footnotesize\color{black!60}\textbf{Ref:} #1}}
\newcommand{\rubrica}[1]{\par\smallskip\noindent\textbf{What earns the points.}\begin{itemize}[leftmargin=1.4em,topsep=1pt]#1\end{itemize}}
\newenvironment{notacao}{\par\smallskip\noindent\textbf{Notation.}\begin{itemize}[leftmargin=1.4em,topsep=1pt,itemsep=0pt]\small}{\end{itemize}}

\begin{document}
\begin{center}
  {\Large\bfseries Macroeconomics I --- MPE Insper}\\[2pt]
  {\large Prova 01 --- detailed solution}\\[2pt]
  \textit{Graded answer: Q1 1.0/1.5, Q2 3.1/5.5, Q3 0.5/3.0}
\end{center}
{\footnotesize\noindent \textbf{How to read this.} Each question and item opens with its statement
reproduced verbatim from the exam, in Portuguese as printed. Each question then has a Notation list. Every
equation a later step uses carries a number, and each step says which equations it uses. Q2 and Q3
keep the exam's own equation labels. Each item ends with the intuition, what the grader penalised
in the submitted answer, and what earns the points. The item-by-item diagnosis is in
\texttt{prova-01-error-analysis.md}.\par}

% =====================================================================
\questao{1}{National accounts (1.5 pt)}
\enunciado{\textbf{1. (1,5 pontos)} Considere uma economia simplificada, sem impostos, gastos do governo e sem
comércio exterior. Nela há dois tipos de firmas: a que produz insumos necessários à produção dos bens,
à cargo da outra firma, vendidos às famílias. Suponha que num determinado ano, a firma produtora de
insumos (Firma A) produziu \$600 que foram vendidos à Firma B, que produziu bens vendidos às famílias
no valor de \$1.000.}

\begin{notacao}
\item A, B, N: the input firm, the final-goods firm and (from 1.b) the social network.
\item $VA_j$: value added of firm $j$, its gross output minus the intermediate inputs it uses up.
\item $C$, $I$, $G$, $NX$: household consumption, investment, government purchases and net exports
(final uses). Here $G=NX=0$.
\item Factor income: wages paid to workers plus profits (operating surplus) left to owners.
\end{notacao}

\passo{0, the three definitions used below}
GDP is the market value of the \emph{final} goods and services produced in the period. It can be
computed three ways, which must agree:
\begin{align}
\text{GDP} &= \textstyle\sum_j VA_j, \tag{G1}\label{G1}\\
VA_j &= \text{gross output}_j-\text{intermediate inputs}_j, \tag{G2}\label{G2}\\
\text{GDP} &= \textstyle\sum_j(\text{wages}_j+\text{profits}_j), \tag{G3}\label{G3}\\
\text{GDP} &= C+I+G+NX. \tag{G4}\label{G4}
\end{align}
A good is \emph{intermediate} when it is used up in producing another good within the period.

\itemq{\textbf{1.a (0,5 ponto)} Calcule o PIB desta economia, explicando precisamente o cálculo e o
conceito utilizado.}
\passo{1, classify}The \$600 of inputs is used up by B within the year, so it is intermediate.
\passo{2, product approach}Using \eqref{G2} for each firm (A buys no inputs):
\begin{align*}
VA_A &= 600-0=600,\\
VA_B &= 1000-600=400.
\end{align*}
Summing with \eqref{G1}:
\[
\text{GDP}=600+400=1000 .
\]
\passo{3, income approach}Each firm's value added is paid out to its workers and owners. A's
factors earn 600 and B's earn 400, so by \eqref{G3}:
\[
\text{GDP}=600+400=1000 .
\]
\passo{4, expenditure approach}The only final use is household consumption, $C=1000$, with
$I=G=NX=0$. By \eqref{G4}:
\[
\text{GDP}=1000+0+0+0=1000 .
\]
\intuicao{Adding $600+1000=1600$ would count the inputs twice: once when A sells them and again
inside the price of B's goods. Netting out intermediate inputs in \eqref{G2} is what makes the three
approaches agree.}
\refbib{Kurlat (2020), §1.1--1.2, pp.~13--19 (the three approaches; Example 1.6, p.~19).}
\rubrica{\item 1000 with the final-good or value-added concept stated (full marks were given).}

\itemq{\textbf{1.b (0,5 ponto)} Considere agora que nesta economia a Firma B contrata o serviço de uma
rede social recentemente criada que é gratuita para seus usuários e financiada por publicidade. A
Firma B gasta \$100 em publicidade na rede social mas isto não gera ainda impacto nas suas vendas no
período, que fica no patamar inicial. Calcule, justificando, o PIB desta economia na ótica do
produto, da renda e do dispêndio, comparando-o com o valor obtido no item anterior (1.a),
explicando o que esta comparação indica.}
\passo{1, classify the advertising}Advertising is a service bought \emph{by a firm} and used up in
its production in the period, so it is an intermediate input of B. The service N gives its users
has a price of zero; no transaction takes place, so it is not in GDP.
\passo{2, product approach}Using \eqref{G2}, with B now using two inputs:
\begin{align*}
VA_A &= 600-0=600,\\
VA_B &= 1000-600-100=300,\\
VA_N &= 100-0=100.
\end{align*}
Summing with \eqref{G1}:
\[
\text{GDP}=600+300+100=1000 .
\]
\passo{3, income approach}A's factors earn 600. B's profit falls by the 100 it pays N, so B's
factors earn 300. N's workers and owners earn 100. By \eqref{G3}:
\[
\text{GDP}=600+300+100=1000 .
\]
\passo{4, expenditure approach}Final uses are unchanged ($C=1000$); the advertising is not a final
use. By \eqref{G4}:
\[
\text{GDP}=1000 .
\]
\begin{center}\small
\begin{tabular}{lcccc}\toprule
Approach & A & B & N & GDP\\\midrule
Product, \eqref{G1}--\eqref{G2} & 600 & 300 & 100 & 1000\\
Income, \eqref{G3} & 600 & 300 & 100 & 1000\\
Expenditure, \eqref{G4} & \multicolumn{3}{c}{$C=1000$} & 1000\\\bottomrule
\end{tabular}\end{center}
\passo{5, comparison}GDP is the same as in 1.a. Only the \emph{distribution} of income moved: 100
went from B's factors to N's. Users do get value from the free service, but GDP does not record it,
because it has no market price. This is a limit of GDP as a welfare measure.
\intuicao{A new firm appears and GDP does not rise, because its whole output is an input into a
product whose final value did not change.}
\erro{The answer gave 1000 with one sentence. The item asked for all three approaches; the grader
wrote ``3 óticas??''.}
\refbib{Kurlat (2020), §1.2 (intermediate vs final); ch.~2 (what GDP leaves out).}
\rubrica{\item Advertising classified as intermediate. \item All three approaches computed and
shown to agree. \item The comparison: same GDP, redistribution, unmeasured value of the free
service.}

\itemq{\textbf{1.c (0,5 ponto)} Suponha agora que a rede social utilize integralmente a receita de
\$100 para produzir um novo software (e nada mais se altere em relação ao que foi descrito em 1.b).
Calcule, justificando, o PIB desta economia na ótica do produto, da renda e do dispêndio,
comparando-o com o valor obtido no item anterior (1.b), explicando o que esta comparação indica.}
\passo{1, classify the software}N pays \$100 to programmers who write software for N's own use. The
software is a durable productive asset, used for more than one period. National accounts record it
as \emph{gross fixed capital formation} (investment in intellectual property), valued at its cost
of 100. It is a final good.
\passo{2, product approach}N's gross output is now ads plus software, $100+100=200$, with no
intermediate inputs. Using \eqref{G2}:
\begin{align*}
VA_N &= 200-0=200,\\
VA_A &= 600,\qquad VA_B=300\ \text{(as in 1.b)}.
\end{align*}
Summing with \eqref{G1}:
\[
\text{GDP}=600+300+200=1100 .
\]
\passo{3, income approach}N pays 100 in wages. Its operating surplus is its output minus wages:
\[
\text{profit}_N=200-100=100 .
\]
So N's factor income is $100+100=200$. By \eqref{G3}:
\[
\text{GDP}=600+300+200=1100 .
\]
\passo{4, expenditure approach}Nothing else changes, so $C=1000$. Investment is now the software,
$I=100$. By \eqref{G4}:
\[
\text{GDP}=1000+100=1100 .
\]
\begin{center}\small
\begin{tabular}{lcccc}\toprule
Approach & A & B & N & GDP\\\midrule
Product, \eqref{G1}--\eqref{G2} & 600 & 300 & 200 & 1100\\
Income, \eqref{G3} & 600 & 300 & wages 100 + profit 100 & 1100\\
Expenditure, \eqref{G4} & \multicolumn{3}{c}{$C=1000$, $I=100$} & 1100\\\bottomrule
\end{tabular}\end{center}
\passo{5, comparison}GDP rises by 100 relative to 1.b because a new \emph{final} good, an investment
good, was produced. If the software were expensed as an intermediate cost, GDP would stay at 1000.
Capitalising software is the modern national-accounts convention, and it is what makes the three
approaches agree at 1100.
\intuicao{In 1.b the \$100 was used up inside B's product. Here it is spent producing something that
lasts, and producing a lasting asset adds to output.}
\erro{The answer reached 1100 but said the \$100 becomes ``wages and consumption''. On the income
side it is wages; on the expenditure side it is \emph{investment}, since consumption is still 1000
by assumption. The three approaches were again missing.}
\refbib{Kurlat (2020), §1.2 (investment as a final use, pp.~17--19).}
\rubrica{\item Software identified as investment, not consumption or intermediate. \item All three
approaches with 1100. \item The comparison: the $+100$ comes from $I$.}

% =====================================================================
\questao{2}{Solow, golden rule and a two-period competitive equilibrium (5.5 pt)}
\enunciado{\textbf{2. (5,5 pontos)} Considere uma economia descrita pelo modelo de Solow, onde as variáveis
agregadas satisfazem as equações: (1) $Y_t=A_t\,.\,(K_t^{\sigma})\,.\,(L_t^{1-\sigma})$,
(2) $K_{t+1}=(1-\delta)K_t+s\,.\,Y_t$, (3) $L_{t+1}=(1+n)L_t$, (4) $A_{t+1}=A_t=A>1$, sendo que
$\{\sigma,\delta,s,n\}\in(0,1)$.}
{\small The statement's equations, numbered as in the exam, for citation below:}
\begin{align}
Y_t &= A_t\,K_t^{\sigma}L_t^{1-\sigma}, \tag{1}\label{e1}\\
K_{t+1} &= (1-\delta)K_t+s\,Y_t, \tag{2}\label{e2}\\
L_{t+1} &= (1+n)L_t, \tag{3}\label{e3}\\
A_{t+1} &= A_t=A>1. \tag{4}\label{e4}
\end{align}
\setcounter{equation}{4}

\begin{notacao}
\item $Y_t$: output (GDP) in period $t$. $K_t$: capital stock. $L_t$: number of workers.
\item $A$: total factor productivity, constant and $>1$ by \eqref{e4}.
\item $\sigma\in(0,1)$: elasticity of output with respect to capital, which under competition is
capital's share of income.
\item $\delta\in(0,1)$: depreciation rate. $s\in(0,1)$: saving rate (share of output saved).
$n\in(0,1)$: growth rate of the labour force.
\item $y_t\equiv Y_t/L_t$, $k_t\equiv K_t/L_t$, $c_t$: output, capital and consumption per worker.
\item $k_{SS}$: steady-state capital per worker. $s_{GR}$: golden-rule saving rate.
\end{notacao}

\itemq{\textbf{2.a (1 ponto)} Defina as variáveis estacionárias por letras minúsculas e encontre uma
expressão para $(k_{t+1}/k_t)$ e, então, o valor de estado estacionário de $k$, denotando-o por
$k_{SS}$.}
\passo{1, define the stationary variables}$A$ is constant and $L_t$ grows at rate $n$, so the
variables that settle down are the per-worker ones:
\begin{equation}
y_t\equiv\frac{Y_t}{L_t},\qquad k_t\equiv\frac{K_t}{L_t}. \label{e5}
\end{equation}
\passo{2, intensive production function}Divide \eqref{e1} by $L_t$, using $A_t=A$ from \eqref{e4}:
\begin{align}
y_t &= \frac{A K_t^{\sigma}L_t^{1-\sigma}}{L_t}\nonumber\\
&= A K_t^{\sigma}L_t^{-\sigma}\nonumber\\
&= A\left(\frac{K_t}{L_t}\right)^{\sigma}\nonumber\\
&= A k_t^{\sigma}. \label{e6}
\end{align}
\passo{3, accumulation per worker}Divide \eqref{e2} by $L_t$:
\[
\frac{K_{t+1}}{L_t}=(1-\delta)\frac{K_t}{L_t}+s\,\frac{Y_t}{L_t}.
\]
Replace $K_t/L_t$ by $k_t$ from \eqref{e5} and $Y_t/L_t$ by $Ak_t^{\sigma}$ from \eqref{e6}:
\begin{align}
\frac{K_{t+1}}{L_t} &= (1-\delta)\,(k_t)+s\,\big(A k_t^{\sigma}\big)\nonumber\\
&= (1-\delta)k_t+s\,A k_t^{\sigma}. \label{e7}
\end{align}
The left side has $K_{t+1}$ over last period's labour. Multiply and divide by $L_{t+1}$, then use
\eqref{e3}:
\begin{align}
\frac{K_{t+1}}{L_t} &= \frac{K_{t+1}}{L_{t+1}}\cdot\frac{L_{t+1}}{L_t}\nonumber\\
&= (1+n)\,k_{t+1}. \label{e8}
\end{align}
Substitute \eqref{e8} into the left side of \eqref{e7}:
\begin{align}
\big((1+n)\,k_{t+1}\big) &= (1-\delta)k_t+s\,A k_t^{\sigma}\nonumber\\
(1+n)\,k_{t+1} &= (1-\delta)k_t+s\,A k_t^{\sigma}. \label{e9}
\end{align}
\passo{4, the ratio}Divide \eqref{e9} by $(1+n)k_t$:
\begin{equation}
\frac{k_{t+1}}{k_t}=\frac{1-\delta+s\,A k_t^{\sigma-1}}{1+n}. \label{e10}
\end{equation}
\passo{5, steady state}In a steady state $k_{t+1}=k_t=k$, so the left side of \eqref{e10} is
$(1)$:
\begin{align*}
(1) &= \frac{1-\delta+s\,A k^{\sigma-1}}{1+n}.
\end{align*}
Multiply by $1+n$ and solve for $k$:
\begin{align*}
1+n &= 1-\delta+s\,A k^{\sigma-1}\\
n+\delta &= s\,A k^{\sigma-1}\\
k^{1-\sigma} &= \frac{sA}{n+\delta}.
\end{align*}
Raise both sides to the power $1/(1-\sigma)$:
\begin{equation}
\boxed{k_{SS}=\left(\frac{sA}{n+\delta}\right)^{\frac{1}{1-\sigma}}} \label{e11}
\end{equation}
\passo{6, stability}In \eqref{e10}, $\sigma<1$ makes $k_t^{\sigma-1}$ fall as $k_t$ rises. So the
ratio is above 1 when $k_t<k_{SS}$ and below 1 when $k_t>k_{SS}$: from any $k_0>0$, $k_t$ converges
to $k_{SS}$.
\intuicao{In the steady state, saving per worker, $sAk^{\sigma}$, exactly covers the capital needed
to equip new workers ($nk$) and to replace worn-out capital ($\delta k$).}
\rubrica{\item \eqref{e6} and the $(1+n)$ factor in \eqref{e8} derived, not assumed. \item The
ratio \eqref{e10} and $k_{SS}$ \eqref{e11} (full marks were given).}

\itemq{\textbf{2.b (1 ponto)} Demonstre qual é a taxa de poupança maximiza o consumo de estado
estacionário (a taxa da chamada regra de ouro, $s_{GR}$).}
\passo{1, consumption per worker}Consumption is the part of output not saved:
\begin{equation}
c_t=(1-s)\,y_t. \label{e12}
\end{equation}
\passo{2, steady-state output}Evaluate \eqref{e6} at $k_{SS}$ from \eqref{e11}:
\begin{align}
y_{SS} &= A\,k_{SS}^{\sigma}\nonumber\\
&= A\left[\left(\frac{sA}{n+\delta}\right)^{\frac{1}{1-\sigma}}\right]^{\sigma}\nonumber\\
&= A\left(\frac{sA}{n+\delta}\right)^{\frac{\sigma}{1-\sigma}}. \label{e13}
\end{align}
The exponent is $\frac{\sigma}{1-\sigma}$: output depends on $k_{SS}$ raised to $\sigma$.
\passo{3, steady-state consumption}Substitute \eqref{e13} for $y$ in \eqref{e12}:
\begin{align}
c_{SS}(s) &= (1-s)\left(A\left(\frac{sA}{n+\delta}\right)^{\frac{\sigma}{1-\sigma}}\right)\nonumber\\
&= (1-s)\,A\left(\frac{sA}{n+\delta}\right)^{\frac{\sigma}{1-\sigma}}. \label{e14}
\end{align}
\passo{4, take logs}Take the log of \eqref{e14}:
\begin{equation}
\ln c_{SS}=\ln A+\ln(1-s)+\frac{\sigma}{1-\sigma}\Big[\ln s+\ln A-\ln(n+\delta)\Big]. \label{e15}
\end{equation}
Only $\ln(1-s)$ and $\frac{\sigma}{1-\sigma}\ln s$ depend on $s$.
\passo{5, first-order condition}Differentiate \eqref{e15} with respect to $s$ and set it to zero:
\begin{equation}
-\frac{1}{1-s}+\frac{\sigma}{(1-\sigma)\,s}=0. \label{e16}
\end{equation}
Multiply \eqref{e16} by $(1-\sigma)s(1-s)>0$:
\begin{align*}
-(1-\sigma)s+\sigma(1-s) &= 0\\
-s+\sigma s+\sigma-\sigma s &= 0\\
\sigma-s &= 0.
\end{align*}
Hence
\begin{equation}
\boxed{s_{GR}=\sigma} \label{e17}
\end{equation}
\passo{6, second-order condition}Differentiate \eqref{e16} once more:
\[
\frac{d^2\ln c_{SS}}{ds^2}=-\frac{1}{(1-s)^2}-\frac{\sigma}{(1-\sigma)s^2}<0 .
\]
So \eqref{e15} is strictly concave in $s$, and \eqref{e17} is the maximum.
\passo{7, cross-check through $k$}In a steady state, \eqref{e9} with $k_{t+1}=k_t=k$ gives
\begin{equation}
sAk^{\sigma}=(n+\delta)\,k. \label{e18}
\end{equation}
Substitute \eqref{e6} for $y$ in \eqref{e12}, then replace $sAk^{\sigma}$ using \eqref{e18}:
\begin{align}
c_{SS} &= (1-s)\big(Ak^{\sigma}\big)\nonumber\\
&= Ak^{\sigma}-sAk^{\sigma}\nonumber\\
&= Ak^{\sigma}-\big((n+\delta)k\big)\nonumber\\
&= Ak^{\sigma}-(n+\delta)k. \label{e19}
\end{align}
Maximise \eqref{e19} over $k$:
\begin{equation}
\sigma A k^{\sigma-1}=n+\delta, \label{e20}
\end{equation}
that is, the marginal product of capital equals $n+\delta$. Solve \eqref{e18} for $s$ and
substitute \eqref{e20} for $n+\delta$:
\begin{align*}
s &= \frac{(n+\delta)\,k}{Ak^{\sigma}}\\
&= \frac{\big(\sigma A k^{\sigma-1}\big)\,k}{Ak^{\sigma}}\\
&= \frac{\sigma A k^{\sigma}}{Ak^{\sigma}}\\
&= \sigma .
\end{align*}
Same answer as \eqref{e17}.
\passo{8, sanity test at a knife edge}If $\sigma\to0$, capital is useless, so the best saving rate
must go to 0. \eqref{e17} gives $s_{GR}\to0$ and passes. The submitted answer $1/(2-\sigma)$ gives
$1/2$ and fails.
\intuicao{Saving more raises $k$, but each extra unit of capital adds less output and needs more
replacement. The golden rule stops where one more unit of capital per worker adds exactly as much
output, $\sigma Ak^{\sigma-1}$, as it costs to maintain, $n+\delta$ (eq.~\eqref{e20}). With
Cobb--Douglas, that happens when the saving rate equals capital's share.}
\erro{The answer wrote $c_{SS}=A(1-s)\left(\frac{sA}{n+\delta}\right)^{\frac{1}{1-\sigma}}$: it
substituted $k_{SS}$ instead of $k_{SS}^{\sigma}$ (step 2 above), so the exponent lost its $\sigma$.
The method (logs, derivative, solving) was right but led to $s=1/(2-\sigma)$.}
\refbib{Kurlat (2020), §4.4 (golden rule); Kurlat ex.~5.8, p.~97.}
\rubrica{\item $c_{SS}$ with exponent $\sigma/(1-\sigma)$, \eqref{e14}. \item FOC solved to
$s_{GR}=\sigma$, with the SOC or the cross-check \eqref{e20}. \item The interpretation: capital
share.}

\itemq{\textbf{2.c (1 ponto)} Considere agora o equilíbrio competitivo de uma economia que existe por
apenas dois períodos. Nesta economia há um estoque inicial de capital que se deprecia totalmente em
um período. E há uma firma representativa que é maximizadora de lucros em mercados perfeitamente
concorrenciais e que tem a mesma função de produção (1). O consumidor representativo escolhe o
quanto consumir e trabalhar no primeiro período, mas atinge idade que o retira da força de trabalho
no segundo período. Logo, suas escolhas dependem de uma riqueza inicial (em termos reais) $a_0>0$ e
da renda real do trabalho apenas no primeiro período: $w_1\,.\,A\,.\,L_1$. A renda não consumida no
primeiro período é transferida para o segundo período através do acúmulo de riqueza, sendo $r$ a
taxa de juros real. Por fim, sua utilidade intertemporal é dada por $U=u(c_1)+\beta u(c_2)$, sendo
$u(c)=ln(c_t)$. Derive e explique a restrição orçamentária intertemporal deste consumidor
representativo.}
\begin{notacao}
\item $c_1$, $c_2$: the household's consumption in periods 1 and 2.
\item $a_0>0$: initial real wealth. $a_1$: assets carried from period 1 to period 2 (saving).
\item $L_1$: hours worked in period 1, with time endowment $\bar L$, normalised to $\bar L=1$.
\item $w_1$: real wage per efficiency unit of labour, so labour income is $w_1AL_1$.
\item $\Pi_1$: profit of the firm, paid to the household that owns it.
\item $r$: real interest rate between periods 1 and 2. $\beta\in(0,1)$: discount factor
(patience).
\item $K_1$: initial capital stock; $k_1=K_1/L_1$, $k_2=K_2/L_1$ (capital per unit of labour).
\end{notacao}
\passo{1, who owns what}The household owns the firm. The firm uses the initial capital stock $K_1$,
which is fully used up in production. After paying wages, the firm's profit
$\Pi_1=Y_1-w_1AL_1$ goes to the household.
\passo{2, period-1 budget}Resources (wealth, wages, profits) equal uses (consumption, saving):
\begin{equation}
c_1+a_1=a_0+w_1AL_1+\Pi_1. \label{e21}
\end{equation}
\passo{3, period-2 budget}The retired household has only its savings with interest, and it consumes
everything because there is no later period:
\begin{equation}
c_2=(1+r)\,a_1. \label{e22}
\end{equation}
Since $c_2\ge0$, \eqref{e22} also implies $a_1\ge0$: nobody lends to a household with no future
income.
\passo{4, combine}Solve \eqref{e22} for $a_1$:
\[
a_1=\frac{c_2}{1+r}.
\]
Substitute this for $a_1$ in \eqref{e21}:
\[
c_1+\left(\frac{c_2}{1+r}\right)=a_0+w_1AL_1+\Pi_1 .
\]
Drop the parentheses:
\begin{equation}
\boxed{c_1+\frac{c_2}{1+r}=a_0+w_1AL_1+\Pi_1} \label{e23}
\end{equation}
together with the time constraint
\begin{equation}
0\le L_1\le 1. \label{e24}
\end{equation}
\passo{5, explain}The left side of \eqref{e23} is the \emph{present value} of lifetime consumption:
period-2 consumption is discounted by $1+r$. The right side is the present value of lifetime
resources (wealth, wages and profits, all received in period 1). Since $u'(c)>0$, leaving
resources unspent is never optimal, so \eqref{e23} holds with equality.
\intuicao{One unit not consumed today becomes $1+r$ units tomorrow. So the price of $c_2$ in units
of $c_1$ is $1/(1+r)$.}
\erro{The answer omitted profits $\Pi_1$ (the grader added ``$+\Pi$ (lucro das firmas)'') and called
the left side ``total consumption adjusted by the interest rate'' instead of its present value.}
\rubrica{\item Both period constraints \eqref{e21}--\eqref{e22}, with profits. \item The combined
constraint \eqref{e23}. \item The present-value interpretation.}

\itemq{\textbf{2.d (1 ponto)} Monte o problema do consumidor representativo descrito em 2.c e calcule
as condições de primeira ordem para \underline{\textbf{todas}} as variáveis de sua escolha (incluindo
a equação descrevendo a decisão intertemporal do consumo). Por fim, encontre a equação que descreve a
escolha ótima de $c_1$ (como função apenas de parâmetros, preços e níveis de renda e riqueza),
interpretando-a economicamente.}
\begin{notacao}
\item $\lambda$: multiplier on the budget \eqref{e23} (marginal utility of a unit of wealth).
\item $\nu\ge0$: multiplier on the time constraint $L_1\le1$ in \eqref{e24}.
\item $W\equiv a_0+w_1A+\Pi_1$: lifetime resources when $L_1=1$ (defined in step 6).
\end{notacao}
\passo{1, list the choices first}The household chooses $c_1$, $c_2$ and $L_1$. Then \eqref{e21}
gives $a_1$. It takes $w_1$, $r$ and $\Pi_1$ as given.
\passo{2, Lagrangian}Maximise $U$ subject to \eqref{e23} and \eqref{e24}:
\begin{equation}
\mathcal L=\ln c_1+\beta\ln c_2
+\lambda\Big[a_0+w_1AL_1+\Pi_1-c_1-\frac{c_2}{1+r}\Big]+\nu\,(1-L_1). \label{e25}
\end{equation}
\passo{3, FOC for $c_1$}Differentiate \eqref{e25} with respect to $c_1$:
\begin{equation}
\frac{1}{c_1}=\lambda. \label{e26}
\end{equation}
\passo{4, FOC for $c_2$}Differentiate \eqref{e25} with respect to $c_2$:
\begin{equation}
\frac{\beta}{c_2}=\frac{\lambda}{1+r}. \label{e27}
\end{equation}
\passo{5, FOC for $L_1$}Differentiate \eqref{e25} with respect to $L_1$:
\begin{equation}
\lambda\,w_1A-\nu=0,\qquad \nu\ge0,\qquad \nu\,(1-L_1)=0. \label{e28}
\end{equation}
Substitute $\lambda$ from \eqref{e26} into \eqref{e28}:
\begin{align*}
\left(\frac{1}{c_1}\right)w_1A-\nu &= 0\\
\nu &= \frac{w_1A}{c_1}>0 .
\end{align*}
Since $\nu>0$, the last condition in \eqref{e28} forces
\begin{equation}
L_1=1. \label{e29}
\end{equation}
Work has no utility cost here (there is no leisure in $U$), while every hour raises income. So the
household supplies its whole time endowment: a corner solution.
\passo{6, Euler equation}Substitute $\lambda$ from \eqref{e26} into \eqref{e27}:
\begin{align*}
\frac{\beta}{c_2} &= \frac{\left(\frac{1}{c_1}\right)}{1+r}\\
\frac{\beta}{c_2} &= \frac{1}{(1+r)\,c_1}.
\end{align*}
Cross-multiply:
\begin{equation}
c_2=\beta(1+r)\,c_1. \label{e30}
\end{equation}
Equivalently, $\frac{1/c_1}{\beta/c_2}=1+r$: the marginal rate of substitution equals the relative
price.
\passo{7, solve}Define lifetime resources at $L_1=1$:
\begin{equation}
W\equiv a_0+w_1A+\Pi_1. \label{e31}
\end{equation}
Substitute \eqref{e30} for $c_2$ and \eqref{e29} for $L_1$ in \eqref{e23}:
\begin{align*}
c_1+\frac{\big(\beta(1+r)\,c_1\big)}{1+r} &= a_0+w_1A\,(1)+\Pi_1 .
\end{align*}
The right side is $W$ by \eqref{e31}; simplify:
\begin{align*}
c_1+\frac{\beta(1+r)\,c_1}{1+r} &= W\\
c_1+\beta c_1 &= W\\
c_1(1+\beta) &= W.
\end{align*}
Hence
\begin{equation}
\boxed{c_1=\frac{a_0+w_1A+\Pi_1}{1+\beta}}. \label{e32}
\end{equation}
Substitute \eqref{e32} for $c_1$ in \eqref{e30}:
\begin{align}
c_2 &= \beta(1+r)\left(\frac{W}{1+\beta}\right)\nonumber\\
&= \frac{\beta(1+r)}{1+\beta}\,W. \label{e33}
\end{align}
From \eqref{e21} with $L_1=1$ \eqref{e29} and $W$ \eqref{e31}, $a_1=W-c_1$. Substitute \eqref{e32} for
$c_1$:
\begin{align}
a_1 &= W-\left(\frac{W}{1+\beta}\right)\nonumber\\
&= \frac{(1+\beta)W-W}{1+\beta}\nonumber\\
&= \frac{\beta}{1+\beta}\,W. \label{e34}
\end{align}
\passo{8, interpret}
\begin{itemize}[leftmargin=1.4em]
\item By \eqref{e32}, the household consumes the fraction $1/(1+\beta)$ of lifetime resources
today: its marginal propensity to consume out of wealth.
\item Differentiating \eqref{e32}: $\partial c_1/\partial\beta=-W/(1+\beta)^2<0$. A more patient
household consumes less today and saves more. If $\beta\to0$, it consumes everything in period 1.
\item \eqref{e32} does not contain $r$. All resources arrive in period 1. A higher $r$ makes saving
more rewarding (substitution effect, $c_1\downarrow$) and makes the saver richer (income effect,
$c_1\uparrow$). With log utility the two cancel exactly.
\end{itemize}
\erro{The Euler equation and $c_1$ were right (full marks on those parts), but there was no FOC for
$L_1$ (``Escolha de $L_1$??'', 0/0.2), and $\Pi_1$ was again missing from $c_1$.}
\refbib{Kurlat (2020), §6.1 (two-period problem, log utility); ch.~9 for households that own firms.}
\rubrica{\item Lagrangian \eqref{e25} with every choice variable. \item FOCs for $c_1$, $c_2$
\emph{and} $L_1$ (corner), \eqref{e26}--\eqref{e29}. \item Euler \eqref{e30}, $c_1$ in closed form
\eqref{e32}, interpretation of $\beta$ and of $r$.}

\itemq{\textbf{2.e (1 ponto)} Calcule a taxa de poupança deste agente representativo (ou seja, da
participação da poupança na renda do agente). Depois encontre uma expressão para $k_2$ como função
apenas de $k_1$ e de parâmetros do modelo (lembrando que se trata de uma economia com mercados
perfeitamente competitivos). Faça então um diagrama da expressão encontrada, colocando no eixo
horizontal $k_1$ e no eixo vertical $k_2$ e mostre \underline{\textbf{graficamente}} qual é o estado
estacionário desta economia}
\begin{notacao}
\item $s^\ast$: the household's saving rate (saving over lifetime resources).
\item $\phi(k_1)$: the law of motion, next period's capital as a function of this period's.
\end{notacao}
\passo{1, saving rate}Divide \eqref{e34} by $W$:
\[
\frac{a_1}{W}=\frac{\left(\frac{\beta}{1+\beta}W\right)}{W}.
\]
Cancel $W$:
\begin{equation}
\boxed{s^\ast\equiv\frac{a_1}{W}=\frac{\beta}{1+\beta}} \label{e35}
\end{equation}
It is constant and set by patience alone.
\passo{2, the firm}With $L_1=1$ from \eqref{e29}, $K_1=k_1$. Using \eqref{e1} with $A_t=A$, the firm
chooses labour to maximise profit:
\begin{equation}
\max_{L_1}\ A k_1^{\sigma}L_1^{1-\sigma}-w_1AL_1. \label{e36}
\end{equation}
Differentiate \eqref{e36} with respect to $L_1$:
\begin{equation}
(1-\sigma)A k_1^{\sigma}L_1^{-\sigma}=w_1A. \label{e37}
\end{equation}
Substitute $L_1=1$ from \eqref{e29} into \eqref{e37} to get the wage bill:
\begin{align}
(1-\sigma)A k_1^{\sigma}\,(1)^{-\sigma} &= w_1A\nonumber\\
w_1A &= (1-\sigma)A k_1^{\sigma}. \label{e38}
\end{align}
\passo{3, profits}Profit is output minus the wage bill, $\Pi_1=Y_1-w_1AL_1$. Substitute $Y_1$ from
\eqref{e1} (with $L_1=1$) and $w_1A$ from \eqref{e38}:
\begin{align}
\Pi_1 &= \big(A k_1^{\sigma}(1)^{1-\sigma}\big)-\big((1-\sigma)A k_1^{\sigma}\big)(1)\nonumber\\
&= A k_1^{\sigma}-(1-\sigma)A k_1^{\sigma}\nonumber\\
&= \sigma A k_1^{\sigma}. \label{e39}
\end{align}
This is capital's share of output. Substitute \eqref{e38} for $w_1A$ and \eqref{e39} for $\Pi_1$ in
\eqref{e31}:
\begin{align}
W &= a_0+\big((1-\sigma)Ak_1^{\sigma}\big)+\big(\sigma Ak_1^{\sigma}\big)\nonumber\\
&= a_0+A k_1^{\sigma}. \label{e40}
\end{align}
\passo{4, market clearing}In a closed economy with no government, the only way to carry resources
into period 2 is to install capital. So all saving becomes period-2 capital:
\begin{equation}
K_2=a_1. \label{e41}
\end{equation}
Substitute \eqref{e34} for $a_1$ in \eqref{e41}, with $k_2=K_2$ (one unit of labour):
\[
k_2=\left(\frac{\beta}{1+\beta}\,W\right).
\]
Then substitute \eqref{e40} for $W$:
\[
k_2=\frac{\beta}{1+\beta}\big(a_0+Ak_1^{\sigma}\big).
\]
Name this map $\phi$:
\begin{equation}
\boxed{k_2=\frac{\beta}{1+\beta}\big(a_0+A k_1^{\sigma}\big)\equiv\phi(k_1)} \label{e42}
\end{equation}
\passo{5, shape of $\phi$}Read off \eqref{e42}:
\begin{itemize}[leftmargin=1.4em]
\item $\phi(0)=\frac{\beta a_0}{1+\beta}>0$: the curve starts above the origin.
\item $\phi'(k)=\frac{\beta}{1+\beta}\,\sigma A k^{\sigma-1}>0$ and $\phi''(k)<0$: increasing and
concave.
\item $\phi'(k)\to\infty$ as $k\to0$ and $\phi'(k)\to0$ as $k\to\infty$.
\end{itemize}
So $\phi(k)-k$ is positive at $k=0$, eventually negative, and strictly concave. It crosses zero
exactly once: one steady state, $\phi(k_{SS})=k_{SS}$. At the crossing the curve cuts the 45° line
from above, so $\phi'(k_{SS})<1$ and the steady state is stable.
\passo{6, steady state}For the economy to repeat, each period's household must face the same
problem. Take later cohorts with no extra wealth ($a_0=0$; the capital they inherit already pays
out through $\Pi$). Substitute $k_2=k_1=k$ and $a_0=0$ into \eqref{e42}:
\begin{align*}
(k) &= \frac{\beta}{1+\beta}\big((0)+A (k)^{\sigma}\big)\\
k &= \frac{\beta}{1+\beta}A k^{\sigma}.
\end{align*}
Divide by $k^{\sigma}$:
\begin{align*}
k^{1-\sigma} &= \frac{\beta A}{1+\beta}.
\end{align*}
Hence
\begin{equation}
\boxed{k_{SS}=\left(\frac{\beta A}{1+\beta}\right)^{\frac{1}{1-\sigma}}} \label{e43}
\end{equation}
The slope there: from step 5,
\[
\phi'(k_{SS})=\sigma\cdot\frac{\beta}{1+\beta}A k_{SS}^{\sigma-1}.
\]
At $k_{SS}$, $k^{1-\sigma}=\frac{\beta A}{1+\beta}$; dividing both sides by $k_{SS}^{1-\sigma}$ gives
$1=\frac{\beta}{1+\beta}Ak_{SS}^{\sigma-1}$. Substitute this bracket into $\phi'(k_{SS})$:
\begin{align*}
\phi'(k_{SS}) &= \sigma\cdot(1)\\
&= \sigma<1 .
\end{align*}
With a recurring $a_0>0$, $k_{SS}$ solves $k=\phi(k)$ in \eqref{e42}; it has no closed form but is
read off the diagram.
\begin{center}
\begin{tikzpicture}
\begin{axis}[width=8.5cm,height=7cm,xmin=0,xmax=0.4,ymin=0,ymax=0.4,
  axis lines=left,xlabel={$k_1$},ylabel={$k_2$},xtick={0.1602,0.2376},xticklabels={$k_{SS}^{a_0=0}$,$k_{SS}$},
  ytick={0.05},yticklabels={$\frac{\beta a_0}{1+\beta}$},tick label style={font=\scriptsize},
  legend style={font=\scriptsize,at={(0.03,0.97)},anchor=north west,draw=none}]
\addplot[black,dashed,domain=0:0.4] {x};
\addlegendentry{45° line}
\addplot[blue!70!black,thick,domain=0.0001:0.4,samples=120] {(0.15+x^0.4)/3};
\addlegendentry{$\phi(k)$, $a_0>0$}
\addplot[red!70!black,thick,dotted,domain=0.0001:0.4,samples=120] {(x^0.4)/3};
\addlegendentry{$\phi(k)$, $a_0=0$}
\addplot[only marks,mark=*,mark size=1.6pt] coordinates {(0.2376,0.2376) (0.1602,0.1602)};
\draw[gray,densely dotted] (axis cs:0.2376,0) -- (axis cs:0.2376,0.2376);
\draw[gray,densely dotted] (axis cs:0.1602,0) -- (axis cs:0.1602,0.1602);
\end{axis}
\end{tikzpicture}\\[-2pt]
{\footnotesize Figure 1. The law of motion \eqref{e42} against the 45° line (illustrative values
$A=1$, $\sigma=0.4$, $\beta=0.5$, $a_0=0.15$). Below $k_{SS}$ the curve lies above the 45° line, so
$k$ rises; above it, $k$ falls.}
\end{center}
\intuicao{Each period the household saves a fixed fraction of what it has \eqref{e35}, and market
clearing \eqref{e41} turns that saving into next period's capital. Diminishing returns ($\sigma<1$)
make output, and so saving, grow less than proportionally with capital. Capital therefore stops
growing where saving just replaces the capital that is used up.}
\par\noindent\textit{Reading note.} The statement does not say who owns $K_1$. In the
overlapping-generations reading (the young save only their wages; the old consume the capital
income), \eqref{e42} becomes $k_2=\frac{\beta}{1+\beta}(1-\sigma)Ak_1^{\sigma}$ and
$k_{SS}=\big(\frac{\beta(1-\sigma)A}{1+\beta}\big)^{1/(1-\sigma)}$. The diagram and the argument are
the same.
\erro{The saving rate was not taken from the optimum (\eqref{e34} already gives it). Instead the
answer set $(1-s)Y=wAL$, which equates consumption with labour income and has no basis in the
model. There was no market-clearing condition \eqref{e41}. The resulting map $k_2=Ak_1^{\sigma}-w_1A$
and the slope inequality were wrong (0/1).}
\refbib{Kurlat (2020), ch.~9 (equilibrium with production); Romer (2012), §2.9--2.10 (Diamond
model, log utility, full depreciation), for the OLG reading.}
\rubrica{\item Saving rate \eqref{e35} from the household's optimum. \item Firm FOC \eqref{e37},
profits \eqref{e39}, and market clearing \eqref{e41}. \item $\phi(k_1)$ \eqref{e42}, the diagram
with the 45° line, and the steady state at the crossing.}

\itemq{\textbf{2.f. (0,5 ponto)} Discuta intuitivamente a relação do modelo dos itens 2.c-2.e com o
modelo de Solow dos itens 2.a-2.b.}
\begin{notacao}
\item $k_{GR}$: the golden-rule capital stock, where \eqref{e20} holds.
\end{notacao}
\passo{1, same law of motion}Substitute $\delta=1$ (capital fully used up) and $n=0$ (constant
labour) into \eqref{e9}:
\begin{align}
\big(1+(0)\big)\,k_{t+1} &= \big(1-(1)\big)k_t+s\,A k_t^{\sigma}\nonumber\\
k_{t+1} &= s\,A k_t^{\sigma}. \label{e44}
\end{align}
With $a_0=0$, \eqref{e42} is exactly \eqref{e44} with $s=\frac{\beta}{1+\beta}$. Likewise,
substitute $\delta=1$, $n=0$ and $s=\frac{\beta}{1+\beta}$ into \eqref{e11}:
\begin{align*}
k_{SS} &= \left(\frac{\left(\frac{\beta}{1+\beta}\right)A}{(0)+(1)}\right)^{\frac{1}{1-\sigma}}\\
&= \left(\frac{\beta A}{1+\beta}\right)^{\frac{1}{1-\sigma}},
\end{align*}
which is \eqref{e43}.
\passo{2, the difference that matters}Solow \emph{assumes} a constant saving rate. Here it is
\emph{chosen}, \eqref{e35}: it depends on patience $\beta$, and it is constant and independent of
$r$ only because utility is logarithmic. With CRRA utility of curvature $\theta\ne1$ it would move
with the interest rate. So the two-period equilibrium is a microfoundation for Solow's saving rule.
\passo{3, against the golden rule}Compare \eqref{e35} with \eqref{e17}. Substitute $\delta=1$ and
$n=0$ into \eqref{e20}:
\begin{align*}
\sigma A k^{\sigma-1} &= (0)+(1)\\
k^{1-\sigma} &= \sigma A\\
k_{GR} &= (\sigma A)^{\frac{1}{1-\sigma}} .
\end{align*}
Comparing with \eqref{e43}:
\[
k_{SS}>k_{GR}\iff\frac{\beta}{1+\beta}>\sigma .
\]
The equilibrium reaches the golden rule only by coincidence. If $\frac{\beta}{1+\beta}>\sigma$, it
accumulates too much capital: saving less would raise consumption in every period (dynamic
inefficiency, possible when households live finitely). Nothing in a competitive equilibrium aims at
maximising steady-state consumption; households weigh today against tomorrow with their own
$\beta$.
\intuicao{Same machine (Solow's accumulation \eqref{e44}), different engine (a saving rate from
preferences rather than a parameter).}
\erro{The answer said only ``Solow with $\delta=1$ and $n=0$'' (0.1/0.5), half of step 1. It missed
the endogenous saving rate and the comparison with $s_{GR}$.}
\refbib{Kurlat (2020), ex.~9.12 ``Optimal vs Fixed Savings Rates'', p.~185.}
\rubrica{\item The mapping $\delta=1$, $n=0$, $s=\beta/(1+\beta)$. \item Endogenous vs exogenous
saving. \item Comparison with $s_{GR}=\sigma$.}

% =====================================================================
\questao{3}{The New Keynesian AS--AD model of Benigno (2015) (3.0 pt)}
\enunciado{\textbf{3. (3 pontos)} Considere o modelo novo-keynesiano simplificado apresentado por Benigno
(2015), que se baseia nas seguintes equações de demanda e oferta agregadas (em que variáveis com
barra em cima são as de longo prazo, quando todas as firmas reajustam preços; as demais sendo de
curto prazo):
\[\text{(AD)}\qquad y=\bar y+(g-\bar g)-\sigma\big[i-(\bar p-p)\big]-\sigma ln\beta\]
\[\text{(AS)}\qquad p-p^e=\kappa(y-y_n)\]
em que $p^e$ é o preço fixado pelas firmas que não podem reajustar preços e $y_n$ é o produto
natural de curto prazo, dado por:
$y_n=\frac{1+\eta}{\sigma^{-1}+\eta}a+\frac{\sigma^{-1}}{\sigma^{-1}+\eta}g-\frac{1}{\sigma^{-1}+\eta}\mu$,
sendo $\mu$ o parâmetro que representa o grau de markup nesta economia.}
{\small The statement's equations, with the exam's labels, for citation below:}
\begin{align}
y &= \bar y+(g-\bar g)-\sigma\big[i-(\bar p-p)\big]-\sigma\ln\beta, \tag{AD}\label{AD}\\
p-p^e &= \kappa\,(y-y_n), \tag{AS}\label{AS}\\
y_n &= \frac{1+\eta}{\sigma^{-1}+\eta}\,a+\frac{\sigma^{-1}}{\sigma^{-1}+\eta}\,g-\frac{1}{\sigma^{-1}+\eta}\,\mu, \tag{YN}\label{YN}\\
L &= \tfrac12(y-y_e)^2+\frac{\phi}{2\kappa}(p-p^e)^2, \tag{LF}\label{LF}\\
y_e &= \frac{1+\eta}{\sigma^{-1}+\eta}\,a+\frac{\sigma^{-1}}{\sigma^{-1}+\eta}\,g. \tag{YE}\label{YE}
\end{align}

\begin{notacao}
\item $y$, $p$: output and price level today (in logs). $\bar y$, $\bar p$: their long-run values,
when all firms have reset prices.
\item $p^e$: the price preset by the firms that cannot adjust this period.
\item $i$: nominal interest rate, the central bank's instrument. $\beta\in(0,1)$: discount factor.
\item $g$, $\bar g$: government spending today and in the long run.
\item $a$, $\bar a$: productivity today and in the long run.
\item $\mu$, $\bar\mu$: the firms' desired markup today and in the long run.
\item $\sigma>0$: intertemporal elasticity of substitution (in Q3; it is \emph{not} the capital
share of Q2). $\eta>0$: inverse Frisch elasticity of labour supply.
\item $\kappa>0$: slope of AS; it is larger the more firms adjust prices.
\item $y_n$: natural output (flexible prices, with the markup). $y_e$: efficient output (no markup).
\item $\phi\ge0$: the weight the central bank puts on price surprises in the loss \eqref{LF}.
\end{notacao}

\passo{0a, long-run output}Long-run output is \eqref{YN} evaluated at long-run values:
\begin{equation}
\bar y=\frac{1+\eta}{\sigma^{-1}+\eta}\,\bar a+\frac{\sigma^{-1}}{\sigma^{-1}+\eta}\,\bar g-\frac{1}{\sigma^{-1}+\eta}\,\bar\mu. \label{e45}
\end{equation}
\textbf{Key fact for Q3:} $\bar\mu$ enters through $\bar y$, which is in \eqref{AD}. The current
$\mu$ enters through $y_n$, which is in \eqref{AS}.
\passo{0b, the two curves in the $(y,p)$ plane}Solve \eqref{AD} for $p$:
\begin{equation}
p=\bar p-i-\ln\beta+\frac{\bar y+g-\bar g-y}{\sigma}. \label{e46}
\end{equation}
This is a downward line with slope $-1/\sigma$. Rewrite \eqref{AS}:
\begin{equation}
p=p^e+\kappa\,(y-y_n). \label{e47}
\end{equation}
This is an upward line with slope $\kappa$ through $(y_n,p^e)$.
\passo{0c, the equilibrium}Substitute \eqref{e47} for $p$ into \eqref{AD}:
\begin{align*}
y &= \bar y+(g-\bar g)-\sigma\Big[i-\Big(\bar p-\big(p^e+\kappa(y-y_n)\big)\Big)\Big]-\sigma\ln\beta .
\end{align*}
Open the brackets:
\begin{align*}
y &= \bar y+g-\bar g-\sigma i+\sigma\bar p-\sigma\big[p^e+\kappa(y-y_n)\big]-\sigma\ln\beta\\
y &= \bar y+g-\bar g-\sigma i+\sigma\bar p-\sigma p^e-\sigma\kappa\,y+\sigma\kappa\,y_n-\sigma\ln\beta .
\end{align*}
Move $-\sigma\kappa\,y$ to the left:
\begin{align*}
y+\sigma\kappa\,y &= \bar y+g-\bar g-\sigma i+\sigma(\bar p-p^e)+\sigma\kappa\,y_n-\sigma\ln\beta .
\end{align*}
Divide by $1+\sigma\kappa$:
\begin{equation}
y=\frac{\bar y+g-\bar g-\sigma i+\sigma(\bar p-p^e)+\sigma\kappa\,y_n-\sigma\ln\beta}{1+\sigma\kappa}. \label{e48}
\end{equation}

\itemq{\textbf{3.a (1 ponto)} Considere que há uma expectativa de aumento no longo prazo do markup.
Represente \underline{graficamente}, explicando economicamente o porquê, a situação de equilíbrio
inicial antes desta informação e o deslocamento que define o novo equilíbrio após a nova
expectativa. Não se esqueça de indicar o hiato do produto e usá-lo em sua explicação.}
\passo{1, initial equilibrium $E_0$}Before the news, monetary policy has set $i$ so that AD crosses
AS at $(y_n,p^e)$: output at its natural level, prices as expected, output gap $y-y_n=0$.
\passo{2, which curve moves}Differentiate \eqref{e45} with respect to $\bar\mu$:
\begin{equation}
d\bar y=-\frac{d\bar\mu}{\sigma^{-1}+\eta}<0. \label{e49}
\end{equation}
By \eqref{YN}, $y_n$ depends on the current $\mu$, which has not moved, so $dy_n=0$.
Hence \eqref{AD} shifts left by $|d\bar y|$ at every $p$ (see \eqref{e46}), and \textbf{AS
\eqref{e47} does not move}. Economically: households expect lower income in the long run; to smooth
consumption (Euler equation) they consume less \emph{today} at any given real interest rate.
\passo{3, new equilibrium $E_1$}Differentiate \eqref{e48}; only $\bar y$, $i$ and $y_n$ can move.
Substitute $di=0$ and $dy_n=0$:
\begin{align}
dy &= \frac{d\bar y-\sigma\,(0)+\sigma\kappa\,(0)}{1+\sigma\kappa}\nonumber\\
&= \frac{d\bar y}{1+\sigma\kappa}. \label{e50}
\end{align}
Substitute \eqref{e49} for $d\bar y$ in \eqref{e50}:
\begin{align*}
dy &= \frac{\left(-\dfrac{d\bar\mu}{\sigma^{-1}+\eta}\right)}{1+\sigma\kappa}\\
&= -\frac{d\bar\mu}{(1+\sigma\kappa)(\sigma^{-1}+\eta)}<0 .
\end{align*}
Differentiate \eqref{e47} ($p^e$ is preset) and substitute $dy_n=0$:
\begin{align}
dp &= \kappa\,\big(dy-(0)\big)\nonumber\\
&= \kappa\,dy<0. \label{e51}
\end{align}
\passo{4, output gap}Since $dy_n=0$, the gap moves one-for-one with output:
\[
d(y-y_n)=dy<0 .
\]
This is a negative output gap, a demand-driven recession, with prices below what the preset firms
expected ($p<p^e$). Flexible-price firms cut prices; firms with preset prices keep $p^e$ and produce
what is demanded, so output falls.
\begin{center}
\begin{tikzpicture}
\begin{axis}[width=9cm,height=6.5cm,xmin=-1,xmax=0.8,ymin=-0.7,ymax=0.6,axis lines=left,
  xlabel={$y$},ylabel={$p$},xtick={-0.4,0},xticklabels={$y_1$,$y_n$},ytick={-0.2,0},yticklabels={$p_1$,$p^e$},
  tick label style={font=\scriptsize},clip=false]
\draw[gray,densely dotted] (axis cs:-1,0) -- (axis cs:0,0) -- (axis cs:0,-0.7);
\draw[gray,densely dotted] (axis cs:-1,-0.2) -- (axis cs:-0.4,-0.2) -- (axis cs:-0.4,-0.7);
\addplot[black,thick,domain=-1:0.8] {0.5*x} node[right,font=\scriptsize]{AS};
\addplot[blue!70!black,thick,domain=-0.55:0.55] {-x} node[below,font=\scriptsize]{AD$_0$};
\addplot[red!70!black,thick,domain=-1:0.1] {-(x+0.6)};
\node[red!70!black,font=\scriptsize,anchor=west] at (axis cs:-0.95,0.4) {AD$_1$};
\addplot[only marks,mark=*,mark size=1.5pt] coordinates {(0,0) (-0.4,-0.2)};
\node[font=\scriptsize,anchor=south west] at (axis cs:0,0) {$E_0=(y_n,p^e)$};
\node[font=\scriptsize,anchor=north east] at (axis cs:-0.4,-0.2) {$E_1$};
\draw[->,red!70!black] (axis cs:0.25,-0.25) -- (axis cs:-0.3,-0.25) node[midway,below,font=\scriptsize]{$\bar\mu\uparrow\Rightarrow\bar y\downarrow$};
\end{axis}
\end{tikzpicture}\\[-2pt]
{\footnotesize Figure 2. News of a higher long-run markup: AD shifts left, AS stays. The output gap is
the horizontal distance $y_1-y_n<0$.}
\end{center}
\intuicao{The markup will raise prices relative to costs only \emph{in the future}. What it changes
today is expected income, which is a demand effect. A supply effect would need firms' costs or
markups to change today.}
\erro{The answer shifted AS through $y_n$ (``$\uparrow\bar\mu$, não $\uparrow\mu$!'') and then
described AD with $p\uparrow,y\downarrow$, a contradictory sign story (0.2/1).}
\refbib{Benigno (2015), §6.3 (news about the future, Fig.~8, p.~513) and §7 (markup shocks, Fig.~9).}
\rubrica{\item AD shifts left because $\bar y$ falls \eqref{e49}; AS fixed because $y_n$ is fixed.
\item Graph with $E_0$ and $E_1$. \item Negative gap \eqref{e50} and the mechanism (consumption
smoothing, sticky prices).}

\itemq{\textbf{3.b (1 ponto)} Com base no item anterior, explique intuitivamente o que uma autoridade
monetária comprometida com a estabilidade de preços deveria fazer após a expectativa de aumento no
markup, e qual o custo econômico de assim proceder.}
\passo{1, the target}Price stability means $p=p^e$. By \eqref{e47}, that requires $y=y_n$, i.e. back
to $E_0$. Monetary policy cannot move AS; it moves AD through $i$.
\passo{2, the size of the move}Differentiate \eqref{e48} and substitute $dy=0$ (the target) and
$dy_n=0$:
\begin{align*}
(0) &= \frac{d\bar y-\sigma\,di+\sigma\kappa\,(0)}{1+\sigma\kappa}\\
0 &= d\bar y-\sigma\,di\\
\sigma\,di &= d\bar y .
\end{align*}
Hence
\begin{equation}
di=\frac{d\bar y}{\sigma}. \label{e52}
\end{equation}
Substitute \eqref{e49} for $d\bar y$ in \eqref{e52}:
\begin{align}
di &= \frac{\left(-\dfrac{d\bar\mu}{\sigma^{-1}+\eta}\right)}{\sigma}\nonumber\\
&= -\frac{d\bar\mu}{\sigma(\sigma^{-1}+\eta)}\nonumber\\
&= -\frac{d\bar\mu}{\sigma\sigma^{-1}+\sigma\eta}\nonumber\\
&= -\frac{d\bar\mu}{1+\sigma\eta}<0. \label{e53}
\end{align}
The central bank \textbf{cuts} $i$. The cut lowers the real rate, raises current consumption, and
moves AD$_1$ back to AD$_0$ in Figure 2. Its size equals the fall in the natural rate of interest.
\passo{3, the cost}
\begin{itemize}[leftmargin=1.4em]
\item In terms of the output gap the cost is \emph{zero}. A pure demand shock has no trade-off:
stabilising prices also closes $y-y_n$ (what the literature calls the divine coincidence).
\item What policy \emph{cannot} fix: output returns only to $y_n$, which is below the efficient
$y_e$ (compare \eqref{YN} with \eqref{YE}: the markup term), and $\bar y$ is permanently lower by
\eqref{e49}. Those are real effects, beyond monetary policy.
\item If $i$ is already near zero, the cut in \eqref{e53} may be infeasible (zero lower bound, §9).
The recession then persists unless policy can promise a higher $\bar p$ or use fiscal policy.
\end{itemize}
\intuicao{Raising $i$, as in the answer, would push AD further left by \eqref{e48}: a deeper
recession and a larger fall in prices, the opposite of price stability.}
\erro{The answer said to \emph{raise} the rate ``to contain the rise in prices''; prices are
falling here. The cost was not discussed (0.1/1).}
\refbib{Benigno (2015), §6 (natural rate and policy), §9 (zero lower bound), §11 (no trade-off for
demand shocks).}
\rubrica{\item Cut $i$, \eqref{e53}, with the reason (AD back to $E_0$). \item No gap cost (divine
coincidence). \item What remains: $y_n<y_e$, lower $\bar y$, the ZLB caveat.}

\itemq{\textbf{3.c (1 ponto)} Suponha que a autoridade monetária representa o conjunto de preferências
sociais e seus objetivos em termos da seguinte função de perda quadrática e do parâmetro $\phi\ge0$:
\[\text{(LF)}\qquad L=\frac12(y-y_e)^2+\frac{\phi}{2\kappa}(p-p^e)^2\]
onde $y_e=\frac{1+\eta}{\sigma^{-1}+\eta}a+\frac{\sigma^{-1}}{\sigma^{-1}+\eta}g$ é o nível eficiente
de produto. Derive a regra ótima de política monetária (IT) e mostre graficamente qual é a política
ótima após a expectativa de aumento do markup no longo prazo. Discuta como sua resposta neste item
depende do valor de $\phi$.}
\begin{notacao}
\item $y^\ast$, $p^\ast$, $i^\ast$: the optimal output, price level and interest rate.
\end{notacao}
\passo{1, the constraint}The central bank sets $i$, which moves only AD. So it can pick any point on
AS, and \eqref{AS} is its constraint:
\begin{equation}
p-p^e=\kappa\,(y-y_n). \label{e54}
\end{equation}
\passo{2, substitute the constraint}Substitute \eqref{e54} for $p-p^e$ in \eqref{LF}:
\begin{align}
L &= \tfrac12(y-y_e)^2+\frac{\phi}{2\kappa}\big(\kappa(y-y_n)\big)^2\nonumber\\
&= \tfrac12(y-y_e)^2+\frac{\phi}{2\kappa}\,\kappa^2(y-y_n)^2\nonumber\\
&= \tfrac12(y-y_e)^2+\frac{\phi\kappa}{2}(y-y_n)^2 . \label{e55}
\end{align}
This is the loss as a function of $y$ alone, $L(y)$.
\passo{3, first-order condition}Differentiate \eqref{e55} and set it to zero:
\begin{equation}
(y-y_e)+\phi\kappa\,(y-y_n)=0. \label{e56}
\end{equation}
The second derivative of \eqref{e55} is $1+\phi\kappa>0$, so \eqref{e56} gives a minimum.
\passo{4, the targeting rule}Group \eqref{e56} as $(y-y_e)+\phi\big(\kappa(y-y_n)\big)=0$. By
\eqref{e54}, the bracket equals $p-p^e$. Substitute it:
\[
(y-y_e)+\phi\big(p-p^e\big)=0 .
\]
That is the targeting rule:
\begin{equation}
\boxed{(y-y_e)+\phi\,(p-p^e)=0}. \label{e57}
\end{equation}
\passo{5, the optimal point}Solve \eqref{e56} for $y$:
\begin{align}
y(1+\phi\kappa) &= y_e+\phi\kappa\,y_n\nonumber\\
y^\ast &= \frac{y_e+\phi\kappa\,y_n}{1+\phi\kappa}. \label{e58}
\end{align}
Subtract $y_n$ from \eqref{e58}:
\begin{align*}
y^\ast-y_n &= \frac{y_e+\phi\kappa\,y_n}{1+\phi\kappa}-y_n\\
&= \frac{y_e+\phi\kappa\,y_n-(1+\phi\kappa)\,y_n}{1+\phi\kappa}\\
&= \frac{y_e-y_n}{1+\phi\kappa}.
\end{align*}
Substitute this for $y-y_n$ in \eqref{e54}:
\begin{align}
p^\ast-p^e &= \kappa\left(\frac{y_e-y_n}{1+\phi\kappa}\right)\nonumber\\
p^\ast &= p^e+\frac{\kappa\,(y_e-y_n)}{1+\phi\kappa}. \label{e59}
\end{align}
\passo{6, the instrument}Substitute $(y^\ast,p^\ast)$ for $(y,p)$ in \eqref{AD}:
\[
(y^\ast)=\bar y+(g-\bar g)-\sigma\big[i-(\bar p-(p^\ast))\big]-\sigma\ln\beta .
\]
Solve for $i$:
\begin{align}
\sigma i &= \bar y+g-\bar g-y^\ast+\sigma(\bar p-p^\ast)-\sigma\ln\beta\nonumber\\
i^\ast &= \frac{\bar y+g-\bar g-y^\ast}{\sigma}+\bar p-p^\ast-\ln\beta. \label{e60}
\end{align}
\passo{7, the news}$\bar\mu$ appears in neither \eqref{YE} nor \eqref{YN}, so by \eqref{e58} and
\eqref{e59} $y^\ast$ and $p^\ast$ do not change. Only $\bar y$ changes. Differentiate \eqref{e60}
and substitute $dy^\ast=dp^\ast=0$:
\begin{align*}
di^\ast &= \frac{d\bar y-(0)}{\sigma}-(0)\\
&= \frac{d\bar y}{\sigma}.
\end{align*}
Substitute \eqref{e49} for $d\bar y$:
\begin{align}
di^\ast &= \frac{\left(-\dfrac{d\bar\mu}{\sigma^{-1}+\eta}\right)}{\sigma}\nonumber\\
&= -\frac{d\bar\mu}{1+\sigma\eta}<0. \label{e61}
\end{align}
$\phi$ enters \eqref{e60} only through $y^\ast$ and $p^\ast$, which did not move, so \eqref{e61}
does not depend on $\phi$. The central bank fully offsets the AD shift for \emph{every} weight: it is
the same cut as \eqref{e53}.
\passo{8, the role of $\phi$}$\phi$ decides \emph{where} on AS the economy sits, not how policy
reacts to this shock. With $\mu>0$, \eqref{YN} and \eqref{YE} give $y_e>y_n$. Then, from
\eqref{e58}--\eqref{e59}:
\begin{itemize}[leftmargin=1.4em]
\item $\phi\to\infty$ (only prices matter): $y^\ast=y_n$, $p^\ast=p^e$, strict price stability.
\item $\phi=0$ (only output matters): $y^\ast=y_e$, $p^\ast>p^e$.
\item $0<\phi<\infty$: $y_n<y^\ast<y_e$; the bank accepts some price surprise to bring output closer
to its efficient level.
\end{itemize}
$\phi$ \emph{would} matter for a \emph{current} markup shock, which moves $y_n$ in \eqref{YN} but
not $y_e$ in \eqref{YE}, and so creates a trade-off.
\begin{center}
\begin{tikzpicture}
\begin{axis}[width=10cm,height=7cm,xmin=-0.3,xmax=1.9,ymin=-0.45,ymax=0.8,axis lines=left,
  xlabel={$y$},ylabel={$p$},xtick={0,0.5,1},xticklabels={$y_n$,$y^\ast$,$y_e$},ytick={0,0.25},
  yticklabels={$p^e$,$p^\ast$},tick label style={font=\scriptsize},clip=false]
\draw[gray,densely dotted] (axis cs:-0.3,0.25) -- (axis cs:0.5,0.25) -- (axis cs:0.5,-0.45);
\draw[gray,densely dotted] (axis cs:-0.3,0) -- (axis cs:1,0) -- (axis cs:1,-0.45);
\addplot[black,thick,domain=-0.3:1.6] {0.5*x} node[right,font=\scriptsize]{AS};
\addplot[gray,dashed,domain=-0.3:1.9] {(1-x)/2} node[below,font=\scriptsize]{targeting rule};
\addplot[blue!70!black,thick,domain=0.05:1.15] {0.75-x} node[below,font=\scriptsize]{AD$_0$};
\addplot[red!70!black,thick,domain=-0.3:0.45] {0.15-x} node[below,font=\scriptsize]{AD$_1$};
\addplot[green!40!black,domain=0:360,samples=120] ({1+0.7071*cos(x)},{0.3536*sin(x)});
\addplot[only marks,mark=*,mark size=1.5pt] coordinates {(0.5,0.25) (0.1,0.05) (1,0)};
\node[font=\scriptsize,anchor=south west] at (axis cs:0.5,0.25) {$O$};
\node[font=\scriptsize,anchor=north east] at (axis cs:0.1,0.05) {$E_1$};
\end{axis}
\end{tikzpicture}\\[-2pt]
{\footnotesize Figure 3. Optimal policy (illustrative values $\kappa=0.5$, $\sigma=1$, $\phi=2$,
$y_n=0$, $y_e=1$, $p^e=0$). The loss ellipse centred at $(y_e,p^e)$ is tangent to AS at $O$, where
the targeting rule \eqref{e57} crosses AS. The news moves AD$_0$ to AD$_1$ (point $E_1$ if $i$ is
unchanged); the optimal cut \eqref{e61} brings it back to AD$_0$, through $O$.}
\end{center}
\intuicao{The loss function says what the bank cares about; AS says what it can get. The news
changes neither, so the best attainable point does not change. Only the interest rate needed to
reach it changes.}
\erro{The answer did not set $dL/dy=0$ along AS (step 3). It substituted AS into AD, differentiated
with respect to $i$, and then used a $y_e$ containing $\mu$, concluding that the response depends on
whether $\phi>\sigma$. For this shock the response \eqref{e61} does not depend on $\phi$ at all.}
\refbib{Benigno (2015), §11, eqs.~(33)--(35), Fig.~14, pp.~521--523; Lista 7 Q4.}
\rubrica{\item Loss written on AS, \eqref{e55}, and $dL/dy=0$ with SOC. \item Targeting rule
\eqref{e57}, $y^\ast$, $p^\ast$, $i^\ast$. \item Graph (ellipse tangent to AS, AD shifted and
restored). \item \eqref{e61} independent of $\phi$, and what $\phi$ does decide.}

\end{document}
